\chapter{Intuitive Proofs}


\begin{principle}[The Pigeonhole Principle]
  \label{pigeonhole}
If $kn+1$ objects are placed into $n$ boxes,
then at least one box has at least $k+1$ objects.
\end{principle}

\begin{example} Given any $101$ integers frrom ${1, 2, 3, ..., 200}$, at least
one of these numbers will divide another.
\end{example}

\begin{scratch}
Since there are $101$ items, we can consider the pigeon hole principle with $k=1$ and $n=100$.

Let us consider the following boxes. Create a box for each of the odd numbers ${1, 3, 5, ..., 199}$ and for any number $x$ if $x$ is of the form $x = 2^k \cdot m$, where $m$ is odd and $k >= 0$, we can put $x$ in the box $m$.

There are $100$ odd numbers in the set so we have $100$ boxes. And any two numbers in a box only differ by $2^k$ for some $k$. Thus, for any two numbers in one box, the smaller number divides the larger one.

For any odd number larger than $101$, it will be the only number in that box.
\end{scratch}

\begin{proof}
  For each number $n$ from the set ${1, 2, 3, ..., 200}$, write it in the form of $n = 2^k \cdot m$ where $k >= 0$ and $m$ is an odd number.

  Now, create a box for each odd number from $1$ to $199$. There will be $100$ such boxes. For each of the given $101$ integers, \\

  If $n = 2^k \cdot m$ then put $n$ in the box numbered $m$. \\

  Since $101$ integers are placed in $100$ boxes, there must be at least one box with more than $1$ integer by $\ref{pigeonhole}$.

  Suppose the box $m$ contains two numbers of the form $n_1 = 2^k \cdot m$ and $n_2 = 2^l \cdot m$ where without loss of generality $k > l$. Then we can show that
  $$\frac{n_1}{n_2} = \frac{2^k \cdot m}{2^l \cdot m} = 2^{k-l}$$

  Here, $2^{k-l}$ is an integer since $k > l$, thus, $n_2$ divides $n_1$.

  Thus, proved.
\end{proof}



